\documentclass[11pt,a4paper]{article}

\usepackage[margin=0.85in]{geometry}
\usepackage[T1]{fontenc}
\usepackage[utf8]{inputenc}
\usepackage{lmodern}
\usepackage{amsmath}
\usepackage{microtype}
\usepackage{booktabs}
\usepackage{array}
\usepackage{graphicx}
\usepackage{enumitem}
\usepackage{xcolor}

\setlist{topsep=0.2em, itemsep=0.2em, parsep=0pt}

\title{\textbf{GabayPoz Recommender System --- Version 1.2}\\[0.4em]
       \large Demonstration of Program-First Recommendations with Q7 Strand Multipliers}
\author{GabayPoz Project --- Team 4 (Model Development)}
\date{\today}

\begin{document}

\maketitle

\begin{abstract}
\noindent
This document is an executable demonstration of the GabayPoz recommender,
version~1.2. It is generated by running \texttt{demo.py} against the bundled
sample fixtures under \texttt{data/sample/} and emitting a LaTeX fragment
(\texttt{results.tex}) that this wrapper includes verbatim. Each subsection
describes one synthetic student profile, including Q7 strand, Q10
affordability, Q11 mobility, Q12 duration tolerance, the top-three program
recommendations, the primary school suggestion, and the rank-1 explanation.
\end{abstract}

\vspace{0.5em}
\hrule
\vspace{0.5em}

\section*{1. About this demonstration}

GabayPoz is a constraint-aware, explainable college program recommender
for Grade~11/12 students in Pozorrubio, Pangasinan. EDA established that
tertiary access is shaped strongly by geography and household economics,
so the recommender keeps Q11 commute and Q10 affordability as hard school
filters. Version~1.2 also restores Q7 as a strand-based aptitude
multiplier, using the tested v1 multiplier map.

The current demo pipeline is:

\begin{enumerate}\setlength\itemsep{0.15em}
  \item \textbf{Student profile.} Q1--Q6 build internal affinity; Q8--Q9
        build aptitude; Q7 multiplies matching aptitude dimensions only.
  \item \textbf{Program ranking.} Programs are ranked first using student
        fit, municipality saturation as small context, and the Q12 penalty.
  \item \textbf{School selection.} For each selected program, feasible
        schools are filtered by Q11 commute and Q10 affordability.
  \item \textbf{Persistence shape.} The model emits exactly three ranked
        rows using the v1.2 model ID; explanations and alternate schools
        remain in the API response.
\end{enumerate}

\noindent
The four profiles below exercise STEM, HUMSS, TVL, and GAS strand inputs.

\vspace{0.5em}
\hrule
\vspace{0.5em}

\section*{2. Demonstrative profiles}

\subsection*{D1: Health-oriented STEM strand, ready for board-exam programs}
\noindent Model: \texttt{tds\_recommender\_v1\_2}; Barangay id: \texttt{1}; Q7: \texttt{STEM}; Q10: \texttt{C}; Q11: \texttt{C}; Q12: \texttt{A}.

\paragraph{Top recommendations.}
\begingroup\small\setlength{\tabcolsep}{6pt}
\noindent\begin{tabular}{@{}c p{4.4cm} p{4.4cm} r@{}}
\toprule
Rank & Program & Primary school & Score \\
\midrule
1 & BS Nursing & PSU Lingayen & 0.746 \\
2 & BS Medical Technology & PSU Lingayen & 0.669 \\
3 & Bachelor of Elementary Education & PSU Lingayen & 0.198 \\
\bottomrule
\end{tabular}\endgroup

\paragraph{Rank-1 explanation.}
BS Nursing matches your strongest Health signals. PSU Lingayen is the primary school suggestion because it passed your travel and affordability limits. Local Health field presence is included only as small context, not a job guarantee.


\subsection*{D2: Arts-oriented HUMSS strand, nearby-only travel}
\noindent Model: \texttt{tds\_recommender\_v1\_2}; Barangay id: \texttt{1}; Q7: \texttt{HUMSS}; Q10: \texttt{B}; Q11: \texttt{A}; Q12: \texttt{B}.

\paragraph{Top recommendations.}
\begingroup\small\setlength{\tabcolsep}{6pt}
\noindent\begin{tabular}{@{}c p{4.4cm} p{4.4cm} r@{}}
\toprule
Rank & Program & Primary school & Score \\
\midrule
1 & AB Communication & PSU Lingayen & 0.760 \\
2 & BS Business Administration & PSU Lingayen & 0.328 \\
3 & BS Computer Science & PSU Lingayen & 0.195 \\
\bottomrule
\end{tabular}\endgroup

\paragraph{Rank-1 explanation.}
AB Communication matches your strongest Arts signals. PSU Lingayen is the primary school suggestion because it passed your travel and affordability limits. Local Arts field presence is included only as small context, not a job guarantee.


\subsection*{D3: Agriculture-oriented TVL strand, moderate budget}
\noindent Model: \texttt{tds\_recommender\_v1\_2}; Barangay id: \texttt{3}; Q7: \texttt{TVL}; Q10: \texttt{B}; Q11: \texttt{B}; Q12: \texttt{B}.

\paragraph{Top recommendations.}
\begingroup\small\setlength{\tabcolsep}{6pt}
\noindent\begin{tabular}{@{}c p{4.4cm} p{4.4cm} r@{}}
\toprule
Rank & Program & Primary school & Score \\
\midrule
1 & BS Agriculture & PSU Lingayen & 0.825 \\
2 & BS Civil Engineering & PSU Lingayen & 0.168 \\
3 & BS Medical Technology & PSU Lingayen & 0.159 \\
\bottomrule
\end{tabular}\endgroup

\paragraph{Rank-1 explanation.}
BS Agriculture matches your strongest Agriculture signals. PSU Lingayen is the primary school suggestion because it passed your travel and affordability limits. Local Agriculture field presence is included only as small context, not a job guarantee.


\subsection*{D4: STEM-oriented GAS strand, moderate budget}
\noindent Model: \texttt{tds\_recommender\_v1\_2}; Barangay id: \texttt{2}; Q7: \texttt{GAS}; Q10: \texttt{B}; Q11: \texttt{B}; Q12: \texttt{B}.

\paragraph{Top recommendations.}
\begingroup\small\setlength{\tabcolsep}{6pt}
\noindent\begin{tabular}{@{}c p{4.4cm} p{4.4cm} r@{}}
\toprule
Rank & Program & Primary school & Score \\
\midrule
1 & BS Computer Science & PSU Lingayen & 0.775 \\
2 & BS Civil Engineering & PSU Lingayen & 0.672 \\
3 & BS Medical Technology & PSU Lingayen & 0.393 \\
\bottomrule
\end{tabular}\endgroup

\paragraph{Rank-1 explanation.}
BS Computer Science matches your strongest STEM signals. PSU Lingayen is the primary school suggestion because it passed your travel and affordability limits. Local STEM field presence is included only as small context, not a job guarantee.


\vspace{1em}
\hrule
\vspace{0.5em}

\section*{3. How to reproduce this PDF}

\begin{verbatim}
cd demo
make            # runs demo.py then pdflatex twice
\end{verbatim}

\noindent
\texttt{make clean} removes the LaTeX intermediates and the generated
\texttt{results.tex}; \texttt{make distclean} additionally removes
\texttt{demo.pdf}. Re-running from scratch reproduces the PDF
deterministically: the recommender is fully stateless, the seven
sample CSVs are version-controlled, and tie-breaking is documented
in TDS \S 5 Phase~5.

\vspace{0.5em}
\hrule
\vspace{0.5em}

\section*{4. References}

\begin{itemize}\setlength\itemsep{0.15em}
  \item \textbf{TDS:} \texttt{docs/tds.md} --- the authoritative
        specification this implementation traces to.
  \item \textbf{EDA v1:} \texttt{docs/eda\_v1.pdf} --- Pangasinan
        education, employment, and access profile.
  \item \textbf{EDA v2:} \texttt{docs/eda\_v2.pdf} --- Pozorrubio
        program access, scholarship analysis, and employment signals.
  \item Burke, R.\ (2002). Hybrid recommender systems.
        \emph{User Modeling and User-Adapted Interaction}, 12(4).
  \item Liu, T.-Y.\ (2009). Learning to rank for information
        retrieval. \emph{Foundations and Trends in IR}, 3(3).
  \item Philippine Institute for Development Studies (2019).
        \emph{Tracer study of CHED-scholarship and non-scholarship
        graduates} (Discussion Paper No.\ 2019-26).
\end{itemize}

\end{document}
